% Preface: the README's "Om kursen", "Efter kursen" and "Att bygga", info/README.md's
% prerequisites and software list, and the course's three verification layers, as the front
% matter of a book.

\chapter{Förord}
\markboth{Förord}{Förord}

Den här boken tar dig från ett grindnät på ett papper till en CAN-nod i en FPGA som en
mikrokontroller kan prata med över SPI.

Första halvan bygger upp VHDL och synkron konstruktion från grunden: kombinatoriska nät,
sekvensnät, metastabilitet och synkronisering, räknare, skiftregister, timers och
tillståndsmaskiner, var och en först resonerad om som en \emph{krets} och först därefter uttryckt
i språket. Andra halvan är ett grupprojekt där den kunskapen används till något större: en
förenklad CAN-kontroller, konstruerad modul för modul, och de två lager som gör den nåbar utifrån
\textemdash{} en registerbank och en SPI-brygga.

\textbf{Ämnet är VHDL och digital konstruktion i FPGA, inte digitalteknikens grunder.} Grindar,
sanningstabeller, boolesk algebra och Karnaughdiagram repeteras snarare än lärs ut, och de
repeteras där språket kräver det. Vad boken lär ut från noll är språket och verktygskedjan.

\section*{Det här förutsätts du ha med dig}
\begin{itemize}
  \item Logiska grindar och sanningstabeller, och att läsa en summa-av-produkter ur en tabell.
  \item Boolesk algebra: de grundläggande identiteterna och De Morgans lagar.
  \item Karnaughdiagram, för funktioner av tre eller fyra variabler.
  \item Register, D-vippor, och vad en klockflank gör.
  \item Tillståndsmaskiner som begrepp: tillstånd, övergångar, och skillnaden mellan Moore och
    Mealy.
\end{itemize}

Inget programmeringsspråk utöver allmän mjukvaruvana förutsätts. Vana vid C eller något liknande
hjälper, eftersom flera förklaringar ställer VHDL:s semantik mot den, och det enda avsnitt som
faktiskt innehåller C++ (\secref{c10:sec:framing}) är en kort inledning till ramning snarare än
ett moment i språket.

Till grupprojektet tillkommer \textbf{Git och GitHub} på grundnivå: klona, branch, commit, push
och Pull Request.

\section*{Det här ska du kunna efteråt}
\begin{itemize}
  \item Realisera en boolesk funktion som ett grindnät, och simulera den för hand.
  \item Läsa och skriva kombinatorisk och sekventiell VHDL: entiteter, arkitekturer, signaler,
    variabler och processer.
  \item Förklara metastabilitet och tillämpa dubbelvippsynkroniseraren på varje asynkron ingång.
  \item Konstruera räknare, skiftregister, timers och tillståndsmaskiner, och komponera en design
    av moduler du redan byggt i stället för att skriva om deras logik.
  \item Konstruera ett kommunikationsprotokolls hårdvara: bittajming, checksummor, bitstoppning,
    ramsekvensering och arbitrering.
  \item Konstruera gränssnittet mellan hårdvara och mjukvara: en registerbank och en
    SPI-transport, mot ett kontrakt någon annan skriver mot.
  \item Verifiera en design genom att köra en självkontrollerande testbänk under GHDL, och spåra
    ett fel tillbaka till koden som orsakade det.
  \item Arbeta i en grupp med feature-branchar, Pull Requests och kodgranskning.
\end{itemize}

\section*{Hur boken är upplagd}
Boken är satt ur en kurs om tjugo pass, och \textbf{varje kapitel är ett pass}. Kapitlets avsnitt
är passets appendix, i samma ordning och med samma underrubriker, så att ett kapitel här och en
föreläsningskatalog i kursrepot går att läsa sida vid sida. De två delarna motsvarar kursens två
halvor:

\begin{booktable}{@{}p{4.6em}p{3.6em}L@{}}
  \thead{Kapitel} & \thead{Pass} & \thead{Ämne} \\
  \midrule
  \multicolumn{3}{@{}l}{\textit{Del \ref{part:vhdl}: digital konstruktion i VHDL}} \\
  1 & \file{L01} & Kombinatorik och första VHDL \\
  2 & \file{L02} & Större nät, multiplexrar och submoduler \\
  3 & \file{L03} & Sekvensnät \\
  4 & \file{L04} & Metastabilitet och synkronisering \\
  5 & \file{L05} & Variabler och hårdvaran under \\
  6 & \file{L06} & Räknare och skiftregister \\
  7 & \file{L07} & Timers \\
  8 & \file{L08} & Tillståndsmaskiner \\
  9 & \file{L09} & Praktisk tentamen 1: sekvensnät \\
  \addlinespace
  \multicolumn{3}{@{}l}{\textit{Del \ref{part:can}: grupprojektet, en CAN-nod}} \\
  10 & \file{L10} & Projektstart: CAN-bussen, ramen och \code{can_def} \\
  11 & \file{L11} & Arkitektur, toppnivån och registerkartan \\
  12 & \file{L12} & Synkronisering och bittimern \\
  13 & \file{L13} & CRC-15-motorn \\
  14 & \file{L14} & Sändningsskiftregister och bitstoppning \\
  15 & \file{L15} & Mottagningsskiftregister och avstoppning \\
  16 & \file{L16} & CAN-kontrollern I: tillståndsmaskinen och sändvägen \\
  17 & \file{L17} & CAN-kontrollern II: mottagning och arbitrering \\
  18 & \file{L18} & Registerbanken och SPI från vågformen \\
  19 & \file{L19} & SPI-bryggan, \code{can_spi_node} och bring-up \\
  20 & \file{L20} & Praktisk tentamen 2: CAN-moduler \\
\end{booktable}

Varje kapitel öppnar med vad som finns i det, arbetar sig igenom materialet, avslutas med en
genomgång av vad du ska kunna förklara efteråt, och slutar med sina övningar. Bilagorna håller
verktygsreferenserna, projektspecifikationen, det delade kontraktet mot drivrutinskursen, och en
avslutande not om var man läser vidare.

Kapitel~\ref{c9:ch} och \ref{c20:ch} är de två examinationstillfällena. De innehåller ingen ny
teori; de säger vad tentamen mäter, vilken form uppgifterna har och hur man förbereder sig, och de
är värda att läsa i förväg snarare än dagen innan.

\textbf{En bok till hör till del~\ref{part:can}.} Den här boken beskriver hur man \emph{bygger} en
CAN-nod; protokollet självt beskrivs i \canbook, en kort svensk bok skriven för båda halvorna av ett
CAN-system. Den är protokollreferensen genom hela andra halvan: dess kapitel 1--6 är
kravspecifikationen för det som byggs här, bit för bit, och dess kapitel~7 beskriver precis den
förenklade kontroller del~\ref{part:can} landar i. Varje kapitel som lägger till en protokollregel
säger var i den boken regeln står, och \secref{app:reading:protocol} har hela kapitelkartan mellan de
två. Läs gärna dess kapitel 1 och 2 innan \chapref{c10:ch}; resten fungerar bäst som uppslagsverk.

\section*{Tre lager av verifiering}
Kursen kontrollerar sitt arbete på tre nivåer, och att hålla isär dem är en poäng i sig:

\begin{booktable}{@{}p{9.2em}L@{}}
  \thead{Lager} & \thead{Vad det kontrollerar} \\
  \midrule
  \textbf{För hand} & I CircuitVerse, innan någon VHDL finns. Du härleder vad kretsen ska göra,
    bygger den, och simulerar den mot din egen förutsägelse. Ingenting automatiskt kontrollerar
    det steget, medvetet: att förutsäga ett beteende och sedan bekräfta det är den färdighet
    resten av boken vilar på. \\
  \textbf{I övningarna} & Du skriver VHDL:en och kör en utdelad självkontrollerande testbänk under
    GHDL. Det behövs inget FPGA-kort, och nästan allt går att verifiera på en vanlig laptop. \\
  \textbf{På kortet} & Designen syntetiseras i Quartus Prime Lite och körs på ett Terasic DE0-CV.
    Under den första halvan demonstreras det steget; i grupprojektets sista pass gör du det
    själv. \\
\end{booktable}

\Secref{c2:sec:testbenches} är den fullständiga introduktionen till testbänkar, och
bilaga~\ref{app:sim} är referensen du kommer tillbaka till. Bilaga~\ref{app:quartus} är
Quartus-flödet.

\section*{Vad övningarna ser ut som}
Övningarna kommer efter varje kapitel, numrerade inom kapitlet, och var och en är märkt med sin
sort:

\begin{booktable}{@{}p{7.4em}L@{}}
  \thead{Sort} & \thead{Vad den ber dig om} \\
  \midrule
  \kindtag[fillaccent5]{Krets} & Rita, bygg och simulera en krets för hand i CircuitVerse, och
    stäm av den mot din egen sanningstabell eller ditt eget tidsdiagram. \\
  \kindtag[fillaccent]{VHDL} & Skriv en modul, med det entitetsnamn och den \textbf{portordning}
    övningen anger, och kör dess utdelade testbänk. \\
  \kindtag[fillaccent3]{Simulering} & Ändra något och ta reda på vad GHDL rapporterar. Flera av
    dem ber dig först \emph{förutsäga} utfallet och sedan köra. \\
  \kindtag[fillaccent2]{Resonemang} & Svara i löpande text eller i siffror: vad koden gör, varför
    en konstruktion är bättre än en annan, eller vad ett tal blir. \\
  \kindtag[fillaccent4]{C++} & Kursens enda C++-moment, i \chapref{c10:ch}. \\
\end{booktable}

\textbf{Det finns inga tryckta lösningar, och det är avsiktligt.} Varje övning som ber dig skriva
en VHDL-modul har en utdelad självkontrollerande testbänk, och den är facit: den driver din modul
genom varje fall och säger exakt vilken kontroll som fallerade. Övningarna med penna och papper
kontrollerar du mot den sanningstabell eller det tidsdiagram du själv härledde, vilket är hela
poängen med dem. En lista med svar skulle ta bort just det steg övningen finns till för.

\begin{aside}
\textbf{Portordningen är kontraktet.} Varje utdelad testbänk binder till modulen den driver
\textbf{positionellt}, och det gör varje instansiering i kursen också. Ingenting binder en port
vid namn. Kontraktet är alltså \emph{ordningen} och \emph{typerna} på portarna, exakt som varje
övnings gränssnittstabell listar dem, uppifrån och ned. Portarnas \emph{namn} är dina.

Den affären är värd att förstå snarare än bara följa. Positionell bindning är kortare att skriva
och att läsa, och den gör gränssnittstabellen till enda sanningskälla. Det den ger upp är
kompilatorns hjälp: byt plats på två portar av samma typ och ingenting klagar vid analys, utan
designen beter sig helt enkelt fel i simuleringen i stället. Boken behåller ändå sina portnamn, så
att texten och din kod talar om samma signaler.
\end{aside}

\section*{Att sätta upp verktygen}
Allt i den här boken körs i en Linuxterminal: på Linux, eller på Windows genom WSL. Tre saker
räcker för allt som inte är ett kort:

\begin{itemize}
  \item \textbf{GHDL}, den öppna VHDL-analysatorn och -simulatorn. Det är verktyget kursen körs
    på: varje modul verifieras mot en testbänk i GHDL, och ingenting du ombeds producera behöver
    något annat.
  \item \textbf{CircuitVerse}, en gratis logiksimulator som körs i webbläsaren
    (\url{https://circuitverse.org/simulator}). Inget att installera.
  \item \textbf{Git}, till grupprojektet, plus \code{g++} till \chapref{c10:ch}:s ramningsövning.
\end{itemize}

\noindent På WSL eller Ubuntu:

\begin{shell}
sudo apt -y update
sudo apt -y install git make ghdl g++
ghdl --version           # 3.x eller senare
\end{shell}

\noindent Kursrepot har testramverket som en submodul, så klona det med submodulerna:

\begin{shell}
git clone --recurse-submodules <repo-url>
git submodule update --init       # om det redan är klonat utan dem
\end{shell}

\noindent Två kommandon i repots rot täcker det mesta boken ber om:

\begin{shell}
make build               # Bygg och simulera allt som finns: exempel, övningar och projekt.
make lint                # Markdownlänkar, dubblettkontroll, inga radslut med blanksteg.
\end{shell}

\noindent\code{make build-project} rapporterar varje projekttestbänk som överhoppad i ett färskt
utcheckat repo. Det är det förväntade tillståndet: \file{controller/} och \file{bridge/}
innehåller bara de utdelade filerna, och modulerna är gruppens att skriva i sitt eget repo. En
testbänk vars moduler saknas hoppas över, aldrig underkänns.

\textbf{Quartus Prime Lite och ett DE0-CV-kort är rekommenderade men inte nödvändiga.} Ingen
övning och inget bedömt moment kräver dem; bilaga~\ref{app:quartus} beskriver flödet för den som
vill, och för grupprojektets sista pass, där varje grupp tar upp sin egen nod på utdelad hårdvara.

\section*{Var man läser vidare}
Den nod som byggs i del~\ref{part:can} är ena halvan av ett system. Den andra halvan, C++-koden
som talar med den över SPI, skrivs i en annan kurs, mot samma registerkarta och samma
protokollspecifikation \textemdash{} och mot samma protokollbok, som är gemensam för de två.
Bilaga~\ref{app:reading} beskriver den, standarderna konstruktionen förhåller sig till, och fyra
riktningar vidare i ämnet.
